% Solution 35.2 — parts (e)–(g)
% Source: samples/01.tex
% Author: Vu Hung Nguyen
% Date: November 2025

\subsection{Solution to Problem 35.2 (parts e--g)}

\begin{solution}
Recall $A_n = \displaystyle\int_{0}^{\frac{\pi}{2}} \cos^{2n} x \, dx$ and $B_n = \displaystyle\int_{0}^{\frac{\pi}{2}} x^2 \cos^{2n} x \, dx$.

\begin{enumerate}
\item[(e)] Summing part (d) from $k = 1$ to $n$ telescopes:
\[
\sum_{k=1}^{n}\frac{1}{k^2}
= 2\left(\frac{B_0}{A_0} - \frac{B_n}{A_n}\right).
\]
Now $A_0 = \dfrac{\pi}{2}$, $B_0 = \dfrac{\pi^3}{24}$, so $\dfrac{B_0}{A_0} = \dfrac{\pi^2}{12}$. Hence
\[
\sum_{k=1}^{n}\frac{1}{k^2} = \frac{\pi^2}{6} - 2\,\frac{B_n}{A_n}.
\]

\item[(f)] On $0 \le x \le \frac{\pi}{2}$, $\sin x \ge \dfrac{2}{\pi}x$, so
\[
\cos x = \sqrt{1 - \sin^2 x} \le \sqrt{1 - \frac{4x^2}{\pi^2}}
\quad\Longrightarrow\quad
\cos^{2n} x \le \left(1 - \frac{4x^2}{\pi^2}\right)^n.
\]
Therefore
\[
B_n \le \int_{0}^{\frac{\pi}{2}} x^2 \left(1 - \frac{4x^2}{\pi^2}\right)^n dx.
\]

\item[(g)] Integrate by parts on the right-hand side of (f) with $u = x$, $dv = x\left(1 - \dfrac{4x^2}{\pi^2}\right)^n dx$.
Since $\displaystyle\int x\left(1 - \frac{4x^2}{\pi^2}\right)^n dx
= -\dfrac{\pi^2}{8(n+1)}\left(1 - \dfrac{4x^2}{\pi^2}\right)^{n+1}$,
the boundary term vanishes and
\[
\int_{0}^{\frac{\pi}{2}} x^2 \left(1 - \frac{4x^2}{\pi^2}\right)^n dx
= \frac{\pi^2}{8(n+1)}\int_{0}^{\frac{\pi}{2}} \left(1 - \frac{4x^2}{\pi^2}\right)^{n+1} dx.
\]
\end{enumerate}
\end{solution}

\begin{takeaways}
\begin{itemize}
    \item Telescoping part (d) expresses the partial sum in terms of $B_n/A_n$
    \item The inequality $\sin x \ge 2x/\pi$ bounds $\cos^{2n} x$ from above
    \item Part (g) evaluates the bounding integral by a further integration by parts
\end{itemize}
\end{takeaways}
